\begin{pastexam}[label={px:ust-115-ee-ds-24}]{臺聯大電機類 115 學年度 第 24 題}
（單選；2 分）

\textbf{If you want to find the top 100 students by scanning a list of scores, which of the
following data structures will offer the best computational time?}
\begin{enumerate}[label=(\Alph*)]
\item min-heap sized 100
\item max-heap sized 100
\item balanced binary search tree such as AVL-tree with 100 nodes
\item sorted linked list sized 100
\item sorted array sized 100
\end{enumerate}
\end{pastexam}

\begin{pastsoln}
答案是 (A)。

設共有 $n$ 筆成績、要前 $k = 100$ 名。 掃描時要一直保留「目前看過的前 $k$ 名」，每讀一筆新成績，
就要和「前 $k$ 名裡最低的那一個」比：比它高就把最低的踢掉、把新的放進去。 所以需要的操作是
「看最小值」和「刪最小值再插入」，正是最小堆積（資結書第~8.3~節）：
\begin{pseudo}{前 $k$ 名}
\State $H \gets$ 前 $k$ 筆成績建成的最小堆積 \Comment{$\Theta(k)$}
\For{其餘每一筆成績 $x$}
    \If{$x > H$ 的最小值} \Comment{看根，$\Theta(1)$}
        \State 把根換成 $x$，再往下沉 \Comment{$O(\log k)$}
    \EndIf
\EndFor
\State \Return $H$ 裡的 $k$ 筆
\end{pseudo}
總時間 $O(n \log k)$，大部分成績只和根比一次就丟掉，實際上接近 $\Theta(n)$。 $k = 100$ 是常數時
就是 $O(n)$。

逐項檢查：
\begin{itemize}
\item (A) 對，見上。
\item (B) 最大堆積的根是前 $100$ 名裡最高的，用不到；要踢掉的最低分在某個葉子，找它要 $\Theta(k)$。 錯。
\item (C) AVL 樹也能做到每筆 $O(\log k)$（找最小、刪除、插入），漸近上和 (A) 一樣，但每次都要
走到最左邊找最小值、插入刪除還要旋轉，常數比堆積大，也多存指標。 題目問「最好」，選 (A)。
\item (D) 排序好的串列：最小值在頭，但插入要找位置，$\Theta(k)$。 錯。
\item (E) 排序好的陣列：可以二分找位置，但插入要搬移，$\Theta(k)$。 錯。
\end{itemize}
\end{pastsoln}
